\chapter[Sermon 81. The Rejoicings and Hymns of Saints Are Recited for Rome Destroyed, and All Ungodliness Taken Away.]{The Rejoicings and Hymns of Saints Are Recited for Rome Destroyed, and All Ungodliness Taken Away.}
\label{sermon:081}
\markright{Sermon 81. The Rejoicings and Hymns of Saints Are Recited for ...}

And after that, I heard the voice of a great multitude in Heaven, saying: Alleluia. Salvation and glory, and honour, and power be ascribed to the Lord our God. For true and righteous are his judgments, because he has judged the great whore, who did corrupt the earth with her fornication, and has avenged the blood of his servants of her hand. And again they said, Alleluia. And the smoke of her ascended forevermore. And the twenty-four elders, and the four beasts fell down, and worshipped God that sat on the throne, saying: Amen. Alleluia. And a voice came out of the throne, saying: Praise our Lord God all those that are his servants, and those that fear him both small and great. And I heard the voice of a great multitude, even as the voice of many waters, and as the voice of great thunderings, saying: Alleluia. For our Lord God omnipotent reigns. Let us be glad and rejoice, and give honour unto him: for the marriage of the Lamb is come.

\index{Antichrist}As the Apostle in this book has abundantly described the oppression of the saints, and the cruel and proud assaults of those who persecute the Gospel—those who mock God and torment his saints—it is said that the complaints of the godly have always arisen, as though God, through his suffering and great patience, were neglecting the oppressed. He also now discourses at length about the rejoicing and praises of the saints, by which they extol the truth and justice of God, never failing to punish the ungodly persecutors. They rejoice here chiefly, and praise God for the removal of Antichrist and all ungodliness with him. This is indeed the first point of this chapter. The second confirms all saints, lest they doubt anything about the salvation of the faithful, which he shows to be most certain. The third recounts the sin of blessed John and the faithful doctrine of the holy Angel, that we should worship no creatures, however holy. Finally, the judge, or avenger, Jesus Christ, coming in judgment is described; moreover, the perdition or punishment of all the ungodly, which the just and holy Lord will take of them, is described. This section, which began in chapter 11 of this book and has been suspended until now, is now finally completed.

\index{Judgment, Last}And indeed the Jubilee of the Saints is diverse, plentiful, and manifold, over the lost and condemned enemies of the godly. First, he hears a voice, and that a great voice of much people in heaven. He shows therefore in general that all heavenly beings—the angels not excepted—sing praises to God in heaven. This we understand shall be at the last judgment, when all the ungodly are trodden under foot. And before these things are done, they are rehearsed and described, so that hereby the godly may comfort themselves in dangers and torments, and may abide steadfast in the true faith, believing that they also, though now oppressed, shall sing praises of thanks to God.

\index{Jerome, Saint}And indeed, he here compiles the entire hymn, spoken in praise of God, the avenger. He places “Alleluia” first, after which he adds the praises, “Salvation and glory,” etc. And “Alleluia” signifies “praise the Lord.” He uses a very common term, well known among all in the primitive church. For certain Psalms have this title, “Halleluyah.” For the chapter thus exhorted and stirred up the people to praise God. So, likewise, now the saints, as if embracing the argument of their song, say, “Alleluia.” These terms have more grace in our and foreign languages than in translation. Thus, they have remained in the church: “O[s]anna,” “Amen,” “Selah,” “Maranatha,” and diverse others. Concerning these, Saint Jerome also writes to Marcella and Damasus.

\index{Rome}Now follows the hymn: salvation and glory, and honor, and so forth. They praise these things in God, and ascribe them wholly to him. This is what I explained in interpreting chapters 4 and 5 of this book. Moreover, they praise God for what is principal in this matter: for his judgments are just and true. This saying seems worthy to be imprinted deeply in the hearts of all men, as it may greatly encourage them in temptations. And because the judgments of God are just and true, he adds, because he has judged the great whore: that is to say, he has taken worthy and fitting punishment of the great whore. Hitherto, the Lord has seemed to many to be too slow and too favorable to Rome and the Roman church, but then they shall see that God is most just. The whore has been spoken of before.

Yet does he repeat here again her most heinous and greatest sins. First, corruption through whoredom and enchantment, whereby is signified seduction by corrupt and wicked doctrine. The later, the shedding of the blood of holy martyrs, whereof we have already spoken many times. Therefore God punishes the corruption of doctrine and the cruelty of the Roman church practiced against the saints of God.

And just as in the beginning they sang “Alleluia,” so also in the end they repeat the same. By this repetition, they declare that the praises which we pour out to God on Earth are most acceptable to God. Immediately following is a sentence that might seem to be from either John or the heavenly dwellers themselves. It signifies that the burning of the ungodly will be perpetual and will never end, as Isaiah said in chapters 30 and 46, and the Lord himself in Matthew 25 and Mark 9. For when he speaks of smoke, he understands that there is fire beneath. Let us earnestly consider these things whenever the pleasures and commodities of Antichrist flatter us. For this perpetual fire is prepared for all the ungodly, especially Antichristians. Then he brings in the twenty-four elders and four beasts praising God, by which the universality of creation is understood. Concerning this, see what is said in chapters 4 and 5 of this book. First, they not only kneel but also fall down, so that we should understand what we ought to do on Earth. They worship God who sits on the throne, neither angels, nor spirits, nor any creatures. Furthermore, with two words he shadows their hymn. For they sing, “Amen” and “Alleluia.” They confirm God to be just, and his judgments to be righteous, and that justly he punishes the whore, and therefore that he is to be praised.

Now also a voice comes out of the throne, namely from God himself, but by the ministry of an angel. For it follows: “Sing praise to our God.” Behold, he says, “our God.” Therefore, he accounts himself here among those who share God with humans. Therefore, he was an angel who recited those things of God. Therefore, what the saints are doing now, they are commanded to do. For in the midst of the praises, this voice is heard from God by the angel. And he commands that they praise, and that the true and only God. He shows moreover who should praise him: all the saints, that is, all who fear God, whether they be great or small. By this commandment, therefore, is signified that God is delighted with the praises of holy men, and approves of them. From this, we now who dwell on Earth, learn to praise the Lord without ceasing, and with a sincere heart. We learn that no one is excepted, of whatever degree or age, sex or condition he be.

Again another hymn is appended, as it were an example of obedience. For God by the angel commands the saints to praise. Now therefore they obey God, and offer to him praises. And how great these praises were, he shows by a double comparison, and by a marvelous brevity, and evident manner. For he says, how the voices of the singers were shrill, as the gushing and noise of many waters: also like the clapping or cracking of great thunders. If such brevity and perspicuity were found in Homer or Virgil, it should have many marvelers thereof, who would extol and commend the elegance. But no man marvels, no man sets forth or commends the elegance and efficacy of the holy Scriptures, lacking example. And again is appended a hymn, the beginning whereof, as of the former, is also Alleluya. And like as in the former hymn the saints have celebrated that God does justly punish the wicked: so in this they preach that God reigns, and shall seem even freely to save the saints. They command therefore to praise the Lord. The reason, for because he is omnipotent, he reigns. He has verily reigned evermore: but since so many things have been permitted by him to the ungodly, many have thought that the ungodly, and chiefly Antichrist, has reigned: but now he has oppressed him, and avenged his glory and his servants, it is made manifest to all men that God alone reigns forevermore. They allege also another cause, why God should be praised, the rather why the godly should be glad and rejoice: for the marriage of the Lamb is come. For as much as that time is now come, wherein the Lamb himself will bring in the children of God, his well-beloved spouse, those I mean whom by his bloodshed he has redeemed, to joys everlasting. Of the marriage shall be spoken a little after at large. Praise and glory be to our Redeemer Christ Jesus the Lord. Amen.
